% SGA 6, Exposición VIII: fragmento español alineado con la fuente.
% Alcance exacto: sección 2 completa, fr.tex, líneas 11290--11451.
% SHA-256 de la fuente francesa completa:
% 0D2158B9FE3B466309A92165727A0B191422E05FBAA0DACD5806E8322C54BA1A.
% Termina inmediatamente antes de la sección 3, sin cierre sintético.
% Las superficies matemáticas se conservan diplomáticamente; las anomalías
% detectadas se notifican por separado para su resolución canónica.

\subsection*{2. Complejo cotangente relativo}

\textbf{2.1.} Daremos aquí una definición del complejo cotangente que, sin
tener toda la generalidad posible (véase Exp. 0 4.4), bastará para las
necesidades del teorema de Riemann--Roch; por supuesto, en el caso considerado,
la definición dada en la Exp. 0 coincide con la que se da aquí.

Sea $f:X\to Y$ un morfismo alisable, es decir, que admite una factorización
\[
\begin{tikzcd}[column sep=large,row sep=large]
X \arrow[r,"i"] \arrow[dr,"f"'] & V \arrow[d,"g"]\\
& Y,
\end{tikzcd}
\]
donde $g$ es un morfismo liso e $i$ una inmersión que se puede suponer cerrada.
Sea $J$ el ideal de $X$ en $V$. La derivación canónica
\[
d:\mathcal O_V\longrightarrow\Omega^1_{V/Y}
\]
define un homomorfismo
\[
J\longrightarrow\mathcal O_V\xrightarrow{d}\Omega^1_{V/Y}\longrightarrow
\Omega^1_{V/Y}\otimes_{\mathcal O_V}\mathcal O_X
\]
que se anula sobre $J^2$ y define, por tanto, un homomorfismo, denotado también
por $d$,
\begin{equation}\tag{2.1.1}
J/J^2\xrightarrow{d} i^*(\Omega^1_{V/Y}).
\end{equation}
Se comprueba fácilmente que este homomorfismo es $\mathcal O_X$-lineal. Se
define entonces un complejo de cadenas $L^{X/Y}_\bullet$ poniendo
\[
L^{X/Y}_0=i^*(\Omega^1_{V/Y}),\qquad
L^{X/Y}_1=J/J^2;
\qquad
L^{X/Y}_i=0\quad\text{si }i\neq0,1,
\]
y tomando como homomorfismo $L^{X/Y}_1\to L^{X/Y}_0$ el homomorfismo (\sisi{1.5.1}{2.1.1}). % Canon SGA-SRC-001511

\textbf{Proposición 2.2.} El complejo $L^{X/Y}_\bullet$ no depende, salvo
isomorfismo canónico en la categoría derivada $D(X)$, de la factorización
elegida por un $Y$-esquema liso.

Sean $V,V'$ dos $Y$-esquemas lisos por los cuales se factoriza $f$. Pongamos
$V''=V\times_Y V'$; entonces $V''$ es liso sobre $Y$, y también sobre $V$ y
$V'$. Además, las inmersiones $X\to V$ y $X\to V'$ definen una inmersión
$X\to V''$ que factoriza $f$. Se puede, pues, reducir la comparación de los
complejos $L^{X/Y}_\bullet$ relativos a dos factorizaciones por $Y$-esquemas
lisos $V$ y $V'$ al caso en que existe un diagrama conmutativo sobre $Y$
\[
\begin{tikzcd}[column sep=large,row sep=large]
& V' \arrow[d,"p"]\\
X \arrow[ur,"j"] \arrow[r,"i"'] & V,
\end{tikzcd}
\]
donde $i$ y $j$ son inmersiones cerradas y $p$ un morfismo liso. Por
funtorialidad, se obtienen entonces homomorfismos
\[
I/I^2\longrightarrow J/J^2;
\qquad
p^*(\Omega^1_{V/Y})\longrightarrow\Omega^1_{V'/Y},
\]
de donde resulta un diagrama
\begin{equation}\tag{2.2.1}
\begin{tikzcd}[column sep=large,row sep=large]
I/I^2 \arrow[r] \arrow[d,"d"'] & J/J^2 \arrow[d,"d"]\\
i^*(\Omega^1_{V/Y}) \arrow[r] & j^*(\Omega^1_{V'/Y}).
\end{tikzcd}
\end{equation}
Como $V$ y $V'$ son lisos sobre $Y$, se obtiene la sucesión exacta
\[
0\longrightarrow p^*(\Omega^1_{V/Y})\longrightarrow\Omega^1_{V'/Y}
\longrightarrow\Omega^1_{V'/V}\longrightarrow0.
\]
Como los módulos de diferenciales considerados son localmente libres sobre
$\mathcal O_{\sisi{V}{V'}}$, al aplicar el funtor $j^*$ se obtiene la sucesión exacta % Canon SGA-SRC-001512
\begin{equation}\tag{2.2.2}
0\longrightarrow i^*(\Omega^1_{V/Y})\longrightarrow j^*(\Omega^1_{V'/Y})
\longrightarrow j^*(\Omega^1_{V'/V})\longrightarrow0.
\end{equation}
Por otra parte, $j$ define una sección $j'$ de $X'=V'\times_V X$:
\[
\begin{tikzcd}[column sep=large,row sep=large]
X' \arrow[r,"i'"] \arrow[d,"p'"'] & V' \arrow[d,"p"]\\
X \arrow[r,"i"'] \arrow[ur,"j"] \arrow[u,"j'"',bend right] & V .
\end{tikzcd}
\]
Sean $I,J,K$ los ideales que definen $i,j,j'$, respectivamente; se tiene
$j=i'\circ j'$, y \sisi{las inmersiones $i',j,j'$ son regulares}{la inmersión $j'$ es regular}; además, $i'$ está % Canon SGA-SRC-001513
definida por el ideal $p^*(I)$ (pues $p$ es plano), y su haz conormal es, por
tanto, $i'^*(p^*(I))=p'^*(I/I^2)$. Se obtiene entonces la sucesión exacta
\begin{equation}\tag{2.2.3}
0\longrightarrow I/I^2\longrightarrow J/J^2\longrightarrow K/K^2\longrightarrow0.
\end{equation}
Demostraremos: a) que el diagrama (2.2.1) es conmutativo; se obtendrá así un
morfismo de complejos $\theta$ del complejo $L^{X/Y}_\bullet$ definido por $V$
al complejo $L^{X/Y}_\bullet$ definido por $V'$; b) que el homomorfismo
$K/K^2\longrightarrow j^*(\Omega^1_{V'/V})$ definido por (2.2.2), (2.2.3) y
(2.2.1) es un isomorfismo; esto implicará que el cilindro de la aplicación
$\theta$ es acíclico, y por tanto que $\theta$ es un cuasi-isomorfismo y define
un isomorfismo en la categoría derivada.

Las propiedades a) y b) son de naturaleza local sobre $X$. Se puede, pues,
suponer que $X,Y,V$ y $V'$ son afines. Pongamos $Y=\Spec(A)$, $V=\Spec(B)$,
$V'=\Spec(B')$, $X=\Spec(C)$, y conservemos las mismas letras para denotar los
homomorfismos y los ideales correspondientes. Se tiene el diagrama
\[
\begin{tikzcd}[column sep=large,row sep=large]
B\otimes_A B \arrow[r,"p\otimes p"] \arrow[d,"\Delta"'] & B'\otimes_A B' \arrow[d,"\Delta"] & B'\otimes_B B' \arrow[l,"\Delta"] \arrow[d]\\
B \arrow[r,"p"] \arrow[dr,"i"'] & B' \arrow[d,"j"] & B'\otimes_B C = B'/IB' \arrow[dl,"p'" ]\\
{} & C & {}
\end{tikzcd}
\]
El homomorfismo $I/I^2\to J/J^2$ se deduce del homomorfismo $IB'\to J$ y
coincide con el que se deduce de $p:I\to J$; el homomorfismo
$i^*(\Omega^1_{V/Y})\longrightarrow j^*(\Omega^1_{V'/Y})$ se deduce del
cuadrado conmutativo (1). Por tanto, la propiedad a) es evidente.

El homomorfismo $K/K^2\to j^*(\Omega^1_{V'/V})$ considerado en b) se puede
definir como sigue: sea $\overline x\in K/K^2$ y sea $x$ un elemento de $J$ que
levanta $\overline x$ (existe porque $K=J/IB'$). Se asocia a $\overline x$ la
imagen de $x$ por el homomorfismo compuesto
$J\xrightarrow{d}\Omega_{B'/A}\otimes_{B'}C\longrightarrow\Omega^1_{B'/B}\otimes_{B'}C$.
Por otra parte, sea $H$ el ideal de aumentación
de $B'\otimes_B B'\longrightarrow B'$; el cuadrado conmutativo (2) da un
homomorfismo $H\to K$, y de él un homomorfismo
$H/H^2\otimes_{B'}C\longrightarrow K/K^2$. Ahora bien,
$H/H^2=\Omega^1_{B'/B}$; por consiguiente, basta comprobar que los
homomorfismos construidos son inversos entre sí, lo cual se hace sin dificultad.

\textbf{Definición 2.3.} Sea $f:X\to Y$ un morfismo alisable. El complejo
$L^{X/Y}_\bullet$, definido salvo isomorfismo canónico en la categoría derivada
$D(X)$, se denomina complejo cotangente relativo de $X$ sobre $Y$ (o de $f$).

\textbf{Proposición 2.4.} Sea $f:X\to Y$ un morfismo de intersección completa
(1.1) alisable. Entonces $L^{X/Y}_\bullet$ es perfecto. Más precisamente, para
toda factorización de $f$ como composición de una inmersión y un morfismo liso,
el complejo $L^{X/Y}_\bullet$ definido en 2.1 es estrictamente perfecto.

Sea $f=g\circ i$, donde $g$ es un morfismo liso e $i$ una inmersión. Por 1.2,
$i$ es una inmersión regular. Si $J$ es el ideal que define $i$, entonces
$J/J^2$ es un $\mathcal O_X$-módulo localmente libre de tipo finito. Del mismo
modo, como $g:V\to Y$ es liso, $i^*(\Omega^1_{V/Y})$ es localmente libre de tipo
finito, de donde resulta 2.4.

\textbf{Corolario 2.5.} Sea $f:X\to Y$ un morfismo de intersección completa
alisable, y sea $f=g\circ i$ una factorización de $f$ por un $Y$-esquema liso
$V$. Sea $J$ el ideal de $\mathcal O_V$ que define $i$. Entonces, si
$K^{\bullet}(X)$ denota el grupo de Grothendieck de los $\mathcal O_X$-módulos
localmente libres de tipo finito (resp. de los complejos perfectos sobre $X$),
el elemento $T_f\in K^{\bullet}(X)$ definido por
\[
T_f=\operatorname{cl}(i^*(\Omega^1_{V/Y}))-\operatorname{cl}(J/J^2)
\]
\[
(\text{resp. }T_f=\operatorname{cl}(L^{X/Y}_\bullet))
\]
no depende de la factorización elegida de $f$. El valor, en un punto $x$ de
$X$, de la aumentación de $T_f$ no es sino la dimensión relativa virtual de $X$
sobre $Y$ en el punto $x$ (1.9).

El elemento $T_f$ está bien definido en los dos casos considerados por 2.4, y
su clase en $K^{\bullet}(X)$ no depende de la factorización: por 2.2 en el
segundo caso, y por la demostración de 2.2 en el primero. La última afirmación
de 2.5 es trivial.

\textbf{Proposición 2.6.} Sean $f:X\to Y$ y $g:Y\to Z$ dos morfismos de
intersección completa alisables, y supongamos además que se satisface la
condición siguiente:

(G) Existe un $Y$-esquema liso $V$ por el cual se factoriza $f$, tal que, para
toda inmersión cerrada $Y\to V'$, se pueda encontrar un $V'$-esquema liso
$V''$, con $V\simeq V''\times_{V'}Y$ (esta condición se satisface, en
particular, si $f$ es cuasiproyectivo y $Y$ posee un haz amplio, tomando como
$V$ un $\mathbb P^r_Y$).

Entonces $g\circ f$ es de intersección completa y alisable, y se tiene un
triángulo distinguido
\[
\begin{tikzcd}[column sep=large,row sep=large]
& L^{X/Y}_\bullet \arrow[dl] \\
Lf^*(L^{Y/Z}_\bullet) \arrow[rr] && L^{X/Z}_\bullet \arrow[ul]
\end{tikzcd}
\]
Sea $V'$ un $Z$-esquema liso por el cual se factoriza $g$. Considérese el
diagrama conmutativo
\[
\begin{tikzcd}[column sep=large,row sep=large]
X \arrow[r,"i"] \arrow[dr,"f"'] & V \arrow[r,"j'"] \arrow[d,"p"] & V'' \arrow[d,"p''"]\\
& Y \arrow[r,"j"] \arrow[dr,"g"'] & V' \arrow[d,"p'"]\\
& & Z .
\end{tikzcd}
\]
La inmersión $j'$ es regular, pues se deduce de $j$ mediante un cambio de base
plano. Por tanto, $j'\circ i$ es una inmersión regular que factoriza
$g\circ f$, el cual es, por consiguiente, alisable y de intersección completa.
Además, se tiene la sucesión exacta de haces conormales
\[
0\longrightarrow i^*(J'/J'^2)\longrightarrow K/K^2\longrightarrow I/I^2\longrightarrow0,
\]
donde $I,J,J',K$ denotan los haces que definen, respectivamente, las
inmersiones $i,j,j',j'\circ i$. Ahora bien, como $p''$ es plano,
$J'=p''^*(J)$. La sucesión puede escribirse, pues,
\begin{equation}\tag{2.6.1}
0\longrightarrow f^*(J/J^2)\longrightarrow K/K^2\longrightarrow I/I^2\longrightarrow0.
\end{equation}
Por otra parte, como $V'$ y $V''$ son lisos sobre $Z$, se tiene la sucesión
exacta
\[
0\longrightarrow p''^*(\Omega^1_{V'/Z})\longrightarrow\Omega^1_{V''/Z}\longrightarrow\Omega^1_{V''/V'}\longrightarrow0.
\]
Aplicando $i^*\circ j'^*$, se obtiene la sucesión exacta
\[
0\longrightarrow (p''\circ j'\circ i)^*(\Omega^1_{V'/Z})\longrightarrow (j'\circ i)^*(\Omega^1_{V''/Z})\longrightarrow (j'\circ i)^*(\Omega^1_{V''/V'})\longrightarrow0.
\]
Como $j'^*(\Omega^1_{V''/V'})=\Omega^1_{V/Y}$ por la hipótesis (G), se obtiene
\begin{equation}\tag{2.6.2}
0\longrightarrow f^*(j^*\Omega^1_{V'/Z})\longrightarrow (j'\circ i)^*(\Omega^1_{V''/Z})\longrightarrow i^*(\Omega^1_{V/Y})\longrightarrow0.
\end{equation}
Combinando (2.6.1) y (2.6.2), basta comprobar que el diagrama
\[
\begin{tikzcd}[column sep=large,row sep=large]
0 \arrow[r] & f^*(J/J^2) \arrow[r] \arrow[d,"f^*(d)"'] & K/K^2 \arrow[r] \arrow[d,"d"] & I/I^2 \arrow[r] \arrow[d,"d"] & 0\\
0 \arrow[r] & f^*(j^*(\Omega^1_{V'/Z})) \arrow[r] & (j'\circ i)^*\Omega^1_{V''/Z} \arrow[r] & i^*\Omega^1_{V/Y} \arrow[r] & 0
\end{tikzcd}
\]
es conmutativo, lo que se verifica localmente sin dificultad volviendo a las
definiciones de los morfismos en juego. De la sucesión exacta de complejos así
obtenida se deduce el triángulo distinguido anunciado en la categoría derivada.

\textbf{Corolario 2.7.} Con las notaciones y las hipótesis de 2.6, se tiene en
$K^{\bullet}(X)$ la igualdad siguiente:
\[
T_{g\circ f}=T_f+f^*(T_g).
\]
